\section{\label{sec:level1}Introduction}

The rapidly evolving field of quantum information processing has generated much interest in quantum optimal control for tasks such as state preparation  \cite{rojan2014arbitrarystate,bondar2023globally}, error correction \cite{waldherr2014quantumerror,gertler2021protecting} and realization of logical gates \cite{scabdelhafez2020universal,scallen2017optimal,schuang2014optimal,scwerninghaus2021leakage}. This methodology has been implemented in various quantum systems \cite{glaser2015training}, such as nuclear magnetic resonance \cite{nmrbretschneider2012conversion,nmrkehlet2004improving,nmrkhaneja2005optimal,nmrsimon2009optimal,nmrspinachhogben2011spinach,nmrxu2008designing}, trapped ions \cite{timuller2015phonon,tinebendahl2009optimal}, nitrogen-vacancy centers in diamonds \cite{nvangerer2013robust,nvchou2015optimal,nvdolde2014high,nvpoggiali2018optimal,nvpoulsen2021optimal,nvrembold2020introduction,nvyan2021population}, Bose-Einstein condensation  \cite{becamri2019optimal,becmennemann2015optimal,becsorensen2018qengine} and neutral atoms \cite{atomguo2019high,atomrosi2013fast}. 

One of the prominent techniques for implementing quantum optimal control is Gradient Ascent Pulse Engineering (GRAPE). The central goal of optimal control is to adjust control parameters in such a way that the infidelity of the desired state transfer or gate operation is minimized. In GRAPE, the minimization is based on gradient descent and thus generally requires the evaluation of gradients. Automatic differentiation (AD) \cite{baydin2018automatic} is a convenient tool that has also made an impact on quantum optimal control with GRAPE \cite{pacleung2017speedup,pacqoc,abdelhafez2019gradient,grapead}.
Internally, AD builds a computational graph and applies the chain rule to automatically compute the gradient of a given function. 

However, the convenience of AD comes at the cost of memory required for storing the full computational graph. The use of semi-automatic differentiation (semi-AD) \cite{semiadgoerz2022quantum} partially reduces the memory overhead compared to full automatic differentiation (full-AD). This reduction can still be insufficient to tackle optimal control for quantum systems of rapidly growing size \cite{arute2019quantum,gong2021quantum}. This motivates us to revisit the original GRAPE scheme \cite{nmrkhaneja2005optimal} which is based on hard-coded gradients (HG), and has the smallest possible memory cost. We perform a scaling analysis of memory usage and runtime, and thereby develop a decision tree guiding the optimal choice of strategy among HG, semi-AD and full-AD. We exemplify the scaling behavior with benchmarks for concrete optimization tasks.

The outline of our paper is as follows. Sec.\ \ref{grapetheory} briefly reviews GRAPE and introduces necessary notation. In Sec.\ \ref{gradients main text}, we illustrate derivation and memory-efficient implementation of the gradients for the most commonly used cost function contributions. Our numerical implementation of HG further improves the efficiency of state propagation. In Sec.\ \ref{Scaling analysis and benchmarking}, we analyze and compare the scaling of runtime and memory usage scaling for HG, AD and semi-AD, and present results from concrete benchmark studies. In Sec.\ \ref{Decision tree}, we formulate a decision tree that facilitates the choice of optimal numerical strategy for GRAPE-based quantum optimal control. We summarize our findings and present our conclusions in Sec.\ \ref{Conclusions}.
